% GENERATED FILE. Edit canonical lessons, then run npm run generate:books.
% canonical-source-hash: fnv1a64:fd803858fbbcce53
% canonical-lessons: JA-C109-kamo, JA-C109-washi, JA-C109-taka, JA-C109-fukurou, JA-C109-kaiko

\chapter{Wild duck, Eagle, Hawk, Owl, Silkworm}
\label{ch:ja-wild-duck-eagle-hawk-owl-silkworm}
\hlchaptermodality{109}

\begin{chapteropening}
\textbf{By the end of this chapter:} \emph{I can say a wild duck, an eagle, a hawk, an owl and a silkworm.}

Complete the last lesson of chapter 109: i can say a wild duck, an eagle, a hawk, an owl and a silkworm.
\end{chapteropening}

\section[\texorpdfstring{\ja{かも} (kamo)}{かも (kamo)}]{\ja{かも} (kamo) — a wild duck}

\label{lesson:JA-C109-kamo}



\begin{quote}
Before the new one: say the Japanese for a squirrel, then the Japanese for a wild boar.
\end{quote}

\subsection*{You'll want to know: \ja{かも}}
\textbf{\ja{かも}} — \emph{kamo} — \textquotedblleft{}a wild duck\textquotedblright{}.

\begin{grammarlens}[title={What you've built}]
One more word for this chapter.
\end{grammarlens}

\subsection*{Your turn}
\begin{itemize}
\raggedright
  \item \emph{Say it:} \emph{kamo}
  \item \emph{Say it:} \emph{kamo}, once more
  \item \emph{Say it:} say \emph{kamo}
  \item \emph{Recall:} say the Japanese for a squirrel, then the Japanese for a wild boar, then say \emph{kamo} again
\end{itemize}

\begin{tcolorbox}[breakable,colback=teal!4,colframe=teal!35!black,title={Before you move on}]
What does \ja{かも} mean? (\textquotedblleft{}A wild duck\textquotedblright{}.) Say it once more.
\end{tcolorbox}
\section[\texorpdfstring{\ja{わし} (washi)}{わし (washi)}]{\ja{わし} (washi) — an eagle}

\label{lesson:JA-C109-washi}



\begin{quote}
Before the new one: say the Japanese for a wild boar, then the Japanese for a wild duck.
\end{quote}

\subsection*{You'll want to know: \ja{わし}}
\textbf{\ja{わし}} — \emph{washi} — \textquotedblleft{}an eagle\textquotedblright{}.

\begin{grammarlens}[title={What you've built}]
One more word for this chapter.
\end{grammarlens}

\subsection*{Your turn}
\begin{itemize}
\raggedright
  \item \emph{Say it:} \emph{washi}
  \item \emph{Say it:} \emph{washi}, once more
  \item \emph{Say it:} say \emph{washi}
  \item \emph{Recall:} say the Japanese for a wild boar, then the Japanese for a wild duck, then say \emph{washi} again
\end{itemize}

\begin{tcolorbox}[breakable,colback=teal!4,colframe=teal!35!black,title={Before you move on}]
What does \ja{わし} mean? (\textquotedblleft{}An eagle\textquotedblright{}.) Say it once more.
\end{tcolorbox}
\section[\texorpdfstring{\ja{たか} (taka)}{たか (taka)}]{\ja{たか} (taka) — a hawk}

\label{lesson:JA-C109-taka}



\begin{quote}
Before the new one: say the Japanese for a wild duck, then the Japanese for an eagle.
\end{quote}

\subsection*{You'll want to know: \ja{たか}}
\textbf{\ja{たか}} — \emph{taka} — \textquotedblleft{}a hawk\textquotedblright{}.

\begin{grammarlens}[title={What you've built}]
One more word for this chapter.
\end{grammarlens}

\subsection*{Your turn}
\begin{itemize}
\raggedright
  \item \emph{Say it:} \emph{taka}
  \item \emph{Say it:} \emph{taka}, once more
  \item \emph{Say it:} say \emph{taka}
  \item \emph{Recall:} say the Japanese for a wild duck, then the Japanese for an eagle, then say \emph{taka} again
\end{itemize}

\begin{tcolorbox}[breakable,colback=teal!4,colframe=teal!35!black,title={Before you move on}]
What does \ja{たか} mean? (\textquotedblleft{}A hawk\textquotedblright{}.) Say it once more.
\end{tcolorbox}
\section[\texorpdfstring{\ja{ふくろう} (fukurou)}{ふくろう (fukurou)}]{\ja{ふくろう} (fukurou) — an owl}

\label{lesson:JA-C109-fukurou}



\begin{quote}
Before the new one: say the Japanese for an eagle, then the Japanese for a hawk.
\end{quote}

\subsection*{You'll want to know: \ja{ふくろう}}
\textbf{\ja{ふくろう}} — \emph{fukurou} — \textquotedblleft{}an owl\textquotedblright{}.

\begin{grammarlens}[title={What you've built}]
One more word for this chapter.
\end{grammarlens}

\subsection*{Your turn}
\begin{itemize}
\raggedright
  \item \emph{Say it:} \emph{fukurou}
  \item \emph{Say it:} \emph{fukurou}, once more
  \item \emph{Say it:} say \emph{fukurou}
  \item \emph{Recall:} say the Japanese for an eagle, then the Japanese for a hawk, then say \emph{fukurou} again
\end{itemize}

\begin{tcolorbox}[breakable,colback=teal!4,colframe=teal!35!black,title={Before you move on}]
What does \ja{ふくろう} mean? (\textquotedblleft{}An owl\textquotedblright{}.) Say it once more.
\end{tcolorbox}
\section[\texorpdfstring{\ja{かいこ} (kaiko)}{かいこ (kaiko)}]{\ja{かいこ} (kaiko) — a silkworm}

\label{lesson:JA-C109-kaiko}



\begin{quote}
Before the new one: say the Japanese for a hawk, then the Japanese for an owl.
\end{quote}

\subsection*{You'll want to know: \ja{かいこ}}
\textbf{\ja{かいこ}} — \emph{kaiko} — \textquotedblleft{}a silkworm\textquotedblright{}.

\begin{grammarlens}[title={What you've built}]
One more word for this chapter.
\end{grammarlens}

\subsection*{Your turn}
\begin{itemize}
\raggedright
  \item \emph{Say it:} \emph{kaiko}
  \item \emph{Say it:} \emph{kaiko}, once more
  \item \emph{Say it:} say \emph{kaiko}
  \item \emph{Recall:} say the Japanese for a hawk, then the Japanese for an owl, then say \emph{kaiko} again
\end{itemize}

\begin{tcolorbox}[breakable,colback=teal!4,colframe=teal!35!black,title={Before you move on}]
What does \ja{かいこ} mean? (\textquotedblleft{}A silkworm\textquotedblright{}.) Say it once more.
\end{tcolorbox}
